\chapter*{Acknowledgments}
\addcontentsline{toc}{chapter}{Acknowledgments}

No work of this scope is created in isolation.

The framework developed in this book emerged from the intersection of several intellectual traditions that long predate it. The mathematical foundations draw from approximation theory, information theory, topology, dynamical systems, and the theory of generalized inverses. The conceptual foundations draw from machine learning, cognitive science, philosophy of science, legal theory, and cosmology. Although the framework assembles these elements into a single structure, its components were developed by many researchers across many disciplines over many decades. I am indebted to all of them, whether or not they appear in the bibliography that follows.

Particular gratitude is owed to the authors whose work provided the starting points for several of the framework's central ideas: the learning-theoretic reconstruction perspective developed by Gunn; influence-function methods introduced by Koh and Liang; the predictive-processing and free-energy traditions associated with Friston and others; the traditions of approximation theory and dynamical systems that supplied much of the mathematical vocabulary; and the broader literature on memory, institutional persistence, and cosmological inference that motivated the framework's applications.

The book also benefited from countless conversations, critiques, objections, and requests for clarification. Many of the ideas presented here became sharper because they survived disagreement. Errors that remain are entirely my own.

I am grateful as well to the communities of independent researchers, open-source contributors, writers, and educators whose willingness to share ideas publicly makes projects of this kind possible. A work that attempts to connect disciplines depends upon the existence of people willing to think beyond disciplinary boundaries.

Finally, I wish to acknowledge a more general debt. Every attempt to understand the past relies on records preserved by others: observations, experiments, archives, books, data sets, and memories that survived long enough to become part of the present. This book studies the conditions under which history remains recoverable after compression. It exists only because so much history was preserved before it.

The framework developed in these pages is offered not as a conclusion but as an invitation. The questions it raises are larger than the answers it provides. Many of its claims remain conditional. Many of its empirical obligations remain open. If the framework proves useful, it will do so because others extend it, challenge it, test it, and discover where it fails.

That is how knowledge advances. Not by preserving everything, but by preserving enough.
